\documentclass[11pt,a4paper]{article}

% ---- Packages ----
\usepackage[utf8]{inputenc}
\usepackage[T1]{fontenc}
\usepackage{amsmath,amssymb,amsthm}
\usepackage{mathtools}
\usepackage[margin=1in,headheight=14pt]{geometry}
\usepackage{hyperref}
\usepackage{enumitem}
\usepackage{xcolor}
\usepackage{tcolorbox}
\usepackage{fancyhdr}
\usepackage{booktabs}

% ---- Hyperref ----
\hypersetup{
    colorlinks=true,
    linkcolor=red,
    citecolor=blue,
    urlcolor=blue
}

% ---- Header ----
\pagestyle{fancy}
\fancyhf{}
\fancyhead[L]{\textit{Lesson 11: Static and Extensive-Form Games}}
\fancyhead[R]{\thepage}
\renewcommand{\headrulewidth}{0.4pt}

% ---- Theorem Environments ----
\theoremstyle{definition}
\newtheorem{definition}{Definition}[section]
\newtheorem{example}{Example}[section]
\newtheorem{exercise}{Exercise}[section]

\theoremstyle{plain}
\newtheorem{theorem}{Theorem}[section]
\newtheorem{lemma}[theorem]{Lemma}
\newtheorem{proposition}[theorem]{Proposition}
\newtheorem{corollary}[theorem]{Corollary}

\theoremstyle{remark}
\newtheorem{remark}{Remark}[section]

% ---- Colored Boxes ----
\tcbuselibrary{theorems,skins,breakable}

\newtcolorbox{keyresult}[1][]{
    colback=green!5!white,
    colframe=green!50!black,
    fonttitle=\bfseries,
    title=#1,
    breakable
}

\newtcolorbox{rlconnection}[1][]{
    colback=blue!5!white,
    colframe=blue!50!black,
    fonttitle=\bfseries,
    title=#1,
    breakable
}

\newtcolorbox{warning}[1][]{
    colback=red!5!white,
    colframe=red!50!black,
    fonttitle=\bfseries,
    title=#1,
    breakable
}

% ---- Operators ----
\newcommand{\R}{\mathbb{R}}
\newcommand{\N}{\mathbb{N}}
\newcommand{\Q}{\mathbb{Q}}
\newcommand{\Z}{\mathbb{Z}}
\newcommand{\C}{\mathbb{C}}
\newcommand{\Prob}{\mathbb{P}}
\newcommand{\E}{\mathbb{E}}
\newcommand{\indicator}{\mathbf{1}}
\DeclareMathOperator{\supp}{supp}
\DeclareMathOperator{\conv}{conv}
\DeclareMathOperator{\argmax}{argmax}
\DeclareMathOperator{\argmin}{argmin}
\DeclareMathOperator{\BR}{BR}

% ---- Title ----
\title{Lesson 33: Normal-Form Games and Pure Nash Equilibrium}
\author{Curriculum: A Rigorous Learning Path to Multi-Agent Deep RL}
\date{}

\begin{document}
\maketitle
\tableofcontents
\newpage

\setcounter{section}{0}

\section{Why Game Theory for RL and ML?}
\label{sec:intro}

In single-agent reinforcement learning, the environment is typically modeled as a
Markov decision process (MDP): the agent chooses actions, the environment responds
according to a fixed (if unknown) transition kernel, and the agent seeks a policy
maximizing cumulative reward. The key simplification is that the environment does
not \emph{adapt} to the agent's behavior.

Multi-agent reinforcement learning (MARL) breaks this assumption. When two or more
learning agents interact in a shared environment, each agent's ``environment'' includes
the other agents, whose behavior changes as they learn. The resulting non-stationarity
means that standard single-agent convergence guarantees fail. To analyze MARL, we need
a mathematical framework for strategic interaction --- and that framework is game theory.

\begin{rlconnection}[MARL Is Game Theory + Learning]
Every MARL algorithm implicitly assumes a game-theoretic solution concept:
\begin{itemize}[nosep]
    \item \textbf{Independent Q-learning}: each agent performs Q-learning ignoring
    others. If it converges, the result is a Nash equilibrium of the stage game
    (under restrictive conditions).
    \item \textbf{Minimax-Q} (Littman, 1994): assumes a two-player zero-sum
    stochastic game and seeks the minimax value --- directly applying
    von Neumann's theorem (Theorem~\ref{thm:minimax}).
    \item \textbf{Nash-Q} (Hu \& Wellman, 2003): computes Nash equilibria of
    stage games at each step.
    \item \textbf{MAPPO, QMIX}: cooperative MARL algorithms whose convergence
    analysis relies on the team game structure (common payoff games).
\end{itemize}
Without understanding Nash equilibria, minimax, and subgame perfection, one cannot
state what these algorithms are trying to compute, let alone prove convergence.
\end{rlconnection}

This lesson covers three progressively richer models of strategic interaction:
\begin{enumerate}[nosep]
    \item \textbf{Normal-form (static) games} (Sections~\ref{sec:normal_form}--\ref{sec:correlated}):
    players choose actions simultaneously.
    \item \textbf{Extensive-form (dynamic) games} (Sections~\ref{sec:extensive}--\ref{sec:spe}):
    players move sequentially, possibly with imperfect information.
    \item \textbf{Solution concepts}: Nash equilibrium, minimax strategies, correlated
    equilibrium, and subgame perfect equilibrium.
\end{enumerate}

\noindent\textbf{Prerequisites.} Familiarity with basic probability (finite probability
spaces suffice for this lesson), linear algebra (convex sets, hyperplane separation),
and real analysis (continuity, compactness in $\R^n$). Knowledge of linear programming
duality is helpful but not required; we state the needed result.

%% ============================================================
%% SECTION 2: NORMAL-FORM GAMES
%% ============================================================

\section{Normal-Form Games}
\label{sec:normal_form}

\subsection{Basic Definitions}
\label{subsec:normal_form_def}

\begin{definition}[Normal-Form Game]
\label{def:normal_form}
A \emph{normal-form game} (or \emph{strategic-form game}) is a tuple
$\Gamma = (N, (S_i)_{i \in N}, (u_i)_{i \in N})$ where:
\begin{enumerate}[nosep,label=(\roman*)]
    \item $N = \{1, 2, \ldots, n\}$ is a finite set of \emph{players}.
    \item For each $i \in N$, $S_i$ is a non-empty set of \emph{strategies}
    (or \emph{actions}) available to player~$i$. We write
    $S = \prod_{i \in N} S_i$ for the set of \emph{strategy profiles}.
    \item For each $i \in N$, $u_i : S \to \R$ is player~$i$'s \emph{payoff function}
    (or \emph{utility function}).
\end{enumerate}
The game is \emph{finite} if each $S_i$ is a finite set.
\end{definition}

\begin{remark}[Notation for ``All but $i$'']
\label{rem:notation}
For a strategy profile $s = (s_1, \ldots, s_n) \in S$, we write
$s_{-i} = (s_1, \ldots, s_{i-1}, s_{i+1}, \ldots, s_n) \in S_{-i}
:= \prod_{j \neq i} S_j$.
We often write $s = (s_i, s_{-i})$ to emphasize player~$i$'s strategy.
\end{remark}

\begin{definition}[Bimatrix Game]
\label{def:bimatrix}
A \emph{bimatrix game} is a two-player finite game presented as a pair of matrices
$(A, B) \in \R^{m \times n} \times \R^{m \times n}$, where player~1 has $m$ strategies,
player~2 has $n$ strategies, $A_{jk} = u_1(j,k)$, and $B_{jk} = u_2(j,k)$.
If $B = -A$, the game is \emph{zero-sum}.
\end{definition}

\subsection{Classical Examples}
\label{subsec:examples}

\begin{example}[Prisoner's Dilemma]
\label{ex:pd}
Two players simultaneously choose to \emph{Cooperate} (C) or \emph{Defect} (D):
\[
\begin{array}{c|cc}
 & C & D \\ \hline
C & (3,3) & (0,5) \\
D & (5,0) & (1,1)
\end{array}
\]
Here $D$ strictly dominates $C$ for both players (Definition~\ref{def:dominance}),
so the unique Nash equilibrium is $(D,D)$ with payoffs $(1,1)$, despite $(C,C)$
giving both players a higher payoff.
\end{example}

\begin{example}[Battle of the Sexes]
\label{ex:bos}
Players~1 and~2 prefer different activities ($O$pera vs.\ $F$ootball) but prefer
coordination to miscoordination:
\[
\begin{array}{c|cc}
 & O & F \\ \hline
O & (3,1) & (0,0) \\
F & (0,0) & (1,3)
\end{array}
\]
This game has two pure Nash equilibria --- $(O,O)$ and $(F,F)$ --- and one mixed
equilibrium (computed in Example~\ref{ex:mixed_bos}).
\end{example}

\begin{example}[Matching Pennies]
\label{ex:mp}
A zero-sum game: each player chooses $H$eads or $T$ails. Player~1 wins if
they match; player~2 wins if they differ:
\[
\begin{array}{c|cc}
 & H & T \\ \hline
H & (1,-1) & (-1,1) \\
T & (-1,1) & (1,-1)
\end{array}
\]
This game has no pure Nash equilibrium (Proposition~\ref{prop:mp_no_pure}),
motivating the extension to mixed strategies.
\end{example}

\begin{example}[Stag Hunt]
\label{ex:stag}
Players choose to hunt a $S$tag (requiring cooperation) or a $H$are (safe but
less rewarding):
\[
\begin{array}{c|cc}
 & S & H \\ \hline
S & (4,4) & (0,3) \\
H & (3,0) & (3,3)
\end{array}
\]
Both $(S,S)$ and $(H,H)$ are pure Nash equilibria. The equilibrium $(S,S)$ is
\emph{payoff-dominant} (higher payoffs for both), while $(H,H)$ is \emph{risk-dominant}
(safer against uncertainty about the opponent).
\end{example}

\begin{rlconnection}[Stag Hunt and Cooperative MARL]
The stag hunt captures a fundamental tension in cooperative MARL: agents can achieve
higher collective reward through coordination, but each agent bears risk if the others
fail to coordinate. This is precisely the challenge faced by CTDE (centralized training,
decentralized execution) algorithms like QMIX and MAPPO.
\end{rlconnection}

\subsection{Dominant and Dominated Strategies}
\label{subsec:dominance}

\begin{definition}[Dominance]
\label{def:dominance}
Let $\Gamma = (N, (S_i), (u_i))$ be a normal-form game and fix player~$i$.
A strategy $s_i \in S_i$ \emph{strictly dominates} $s_i' \in S_i$ if
\[
    u_i(s_i, s_{-i}) > u_i(s_i', s_{-i}) \quad \text{for all } s_{-i} \in S_{-i}.
\]
A strategy $s_i$ \emph{weakly dominates} $s_i'$ if
\[
    u_i(s_i, s_{-i}) \geq u_i(s_i', s_{-i}) \quad \text{for all } s_{-i} \in S_{-i},
\]
with strict inequality for at least one $s_{-i}$.
A strategy is \emph{strictly dominated} if some other strategy strictly dominates it.
A strategy is a \emph{(strictly) dominant strategy} if it strictly dominates every
other strategy of that player.
\end{definition}

\begin{proposition}[Dominated Strategies Are Never Best Responses]
\label{prop:dominated_not_br}
If $s_i'$ is strictly dominated by $s_i$, then $s_i'$ is never a best response:
for every $s_{-i} \in S_{-i}$, $s_i' \notin \argmax_{t_i \in S_i} u_i(t_i, s_{-i})$.
\end{proposition}

\begin{proof}
Fix any $s_{-i} \in S_{-i}$. Since $s_i$ strictly dominates $s_i'$, we have
$u_i(s_i, s_{-i}) > u_i(s_i', s_{-i})$, so $s_i'$ does not maximize $u_i(\cdot, s_{-i})$.
\end{proof}

\subsection{Iterated Elimination of Strictly Dominated Strategies}
\label{subsec:iesds}

\begin{definition}[IESDS]
\label{def:iesds}
\emph{Iterated elimination of strictly dominated strategies} (IESDS) is the following
procedure. Set $S_i^0 = S_i$ for all~$i$. For $k = 0, 1, 2, \ldots$, define
\[
    S_i^{k+1} = \big\{s_i \in S_i^k : \text{there is no } s_i' \in S_i^k
    \text{ such that } u_i(s_i', s_{-i}) > u_i(s_i, s_{-i})
    \text{ for all } s_{-i} \in S_{-i}^k\big\}.
\]
The set of \emph{surviving strategies} for player~$i$ is
$S_i^\infty = \bigcap_{k=0}^\infty S_i^k$.
\end{definition}

\begin{theorem}[Order Independence of IESDS]
\label{thm:iesds_order}
In a finite game, the set of strategy profiles surviving iterated elimination of
strictly dominated strategies is independent of the order of elimination.
\end{theorem}

\begin{proof}
We prove that any strategy eliminated in some order of elimination is eventually
eliminated in every order.

Let $s_i^*$ be a strategy of player~$i$ that is eliminated at round~$k$ in some
particular elimination order $E$. This means there exists $s_i' \in S_i^{k-1}(E)$
such that $u_i(s_i', s_{-i}) > u_i(s_i^*, s_{-i})$ for all $s_{-i} \in S_{-i}^{k-1}(E)$.

Consider an alternative elimination order $E'$. We show $s_i^*$ is eventually
eliminated under $E'$. Let $T = S_{-i}^\infty(E')$ be the surviving opponents'
strategies under $E'$.

Since $S_{-i}^\infty(E') \subseteq S_{-i}^{k-1}(E')$ and every strategy surviving
under any order is in particular not eliminated in the first $k-1$ rounds under $E$
(we prove this by induction on rounds), we need a cleaner argument.

We proceed by induction on the round of elimination. For the base case ($k=0$),
$s_i^*$ is strictly dominated in the original game, so it will be eliminated in
round~0 under any order.

For the inductive step, suppose every strategy eliminated in rounds $0, \ldots, k-1$
under \emph{any} elimination order is eventually eliminated under every order. Now
suppose $s_i^*$ is eliminated at round~$k$ under order~$E$. Then there exists
$s_i' \in S_i^{k-1}(E)$ with $u_i(s_i', s_{-i}) > u_i(s_i^*, s_{-i})$ for all
$s_{-i} \in S_{-i}^{k-1}(E)$.

Under order $E'$, the set $S_{-i}^{k-1}(E)$ may differ from $S_{-i}^{k-1}(E')$.
However, the key observation is:

\emph{Claim}: $S_{-i}^\infty(E') \subseteq S_{-i}^{k-1}(E)$.

To prove the claim, note that any strategy $s_j$ removed before round~$k$ under
order~$E$ was strictly dominated by some $s_j'$ against strategies surviving at
that point. By the inductive hypothesis, $s_j$ is eventually removed under $E'$ as
well, so $s_j \notin S_j^\infty(E')$. Taking contrapositives:
$s_j \in S_j^\infty(E') \implies s_j \in S_j^{k-1}(E)$, establishing the claim.

Now, since $S_{-i}^\infty(E') \subseteq S_{-i}^{k-1}(E)$, the domination
$u_i(s_i', s_{-i}) > u_i(s_i^*, s_{-i})$ holds for all $s_{-i} \in S_{-i}^\infty(E')$.
If $s_i' \in S_i^\infty(E')$, then $s_i^*$ is dominated (by $s_i'$) relative to
the surviving strategies under $E'$ and hence is eliminated. If $s_i'$ is itself
eliminated under $E'$, then $s_i'$ was dominated by some $s_i''$, and by transitivity
of strict domination on the surviving set, $s_i^*$ remains dominated and is eliminated.

By induction, the surviving sets coincide for all elimination orders.
\end{proof}

\subsection{Rationalizability}
\label{subsec:rationalizability}

\begin{definition}[Best Response]
\label{def:best_response}
Player~$i$'s \emph{best response correspondence} is the set-valued map
$\BR_i : S_{-i} \rightrightarrows S_i$ defined by
\[
    \BR_i(s_{-i}) = \argmax_{s_i \in S_i} u_i(s_i, s_{-i}).
\]
\end{definition}

\begin{definition}[Rationalizable Strategies]
\label{def:rationalizable}
A strategy $s_i$ is \emph{rationalizable} if it survives the iterated elimination
of strategies that are never best responses. Formally, set $R_i^0 = S_i$ and define
\[
    R_i^{k+1} = \big\{s_i \in R_i^k : \exists\, s_{-i} \in R_{-i}^k
    \text{ such that } s_i \in \BR_i(s_{-i})\big\}.
\]
The set of rationalizable strategies is $R_i^\infty = \bigcap_{k=0}^\infty R_i^k$.
\end{definition}

\begin{proposition}[IESDS and Rationalizability in Two-Player Games]
\label{prop:iesds_rational}
In finite two-player games, the set of rationalizable strategies coincides with the
set of strategies surviving IESDS.
\end{proposition}

\begin{proof}
Every strictly dominated strategy is never a best response
(Proposition~\ref{prop:dominated_not_br}), so
$S_i^\infty \supseteq R_i^\infty$. Conversely, in a two-player finite game,
a strategy that is never a best response to any pure strategy of the opponent is
strictly dominated (this follows from the supporting hyperplane theorem applied
to the payoff polyhedron). Hence a strategy eliminated in the rationalizability
procedure is strictly dominated relative to the surviving set, giving
$R_i^\infty \supseteq S_i^\infty$.
\end{proof}

\begin{warning}[Three or More Players]
For games with three or more players, IESDS may not eliminate all never-best-response
strategies. Rationalizability can be a strictly coarser concept than IESDS survival
in games with $n \geq 3$ players when we restrict to pure strategies. The equivalence
is restored when we extend dominance to allow domination by mixed strategies.
\end{warning}

%% ============================================================
%% SECTION 3: NASH EQUILIBRIUM
%% ============================================================

\section{Nash Equilibrium in Pure Strategies}
\label{sec:nash_pure}

\begin{definition}[Pure Strategy Nash Equilibrium]
\label{def:nash_pure}
A strategy profile $s^* = (s_1^*, \ldots, s_n^*) \in S$ is a \emph{(pure strategy)
Nash equilibrium} (NE) if no player can unilaterally improve their payoff:
\[
    u_i(s_i^*, s_{-i}^*) \geq u_i(s_i, s_{-i}^*)
    \quad \text{for all } s_i \in S_i, \text{ for all } i \in N.
\]
Equivalently, $s_i^* \in \BR_i(s_{-i}^*)$ for all $i \in N$.
\end{definition}

\begin{proposition}[NE as Mutual Best Response]
\label{prop:ne_br}
A strategy profile $s^*$ is a Nash equilibrium if and only if
$s_i^* \in \BR_i(s_{-i}^*)$ for every $i \in N$.
\end{proposition}

\begin{proof}
By Definition~\ref{def:nash_pure}, $s^*$ is a NE if and only if
$u_i(s_i^*, s_{-i}^*) \geq u_i(s_i, s_{-i}^*)$ for all $s_i \in S_i$ and all~$i$.
By Definition~\ref{def:best_response}, this is precisely
$s_i^* \in \argmax_{s_i \in S_i} u_i(s_i, s_{-i}^*) = \BR_i(s_{-i}^*)$ for all~$i$.
\end{proof}

\begin{proposition}[Matching Pennies Has No Pure NE]
\label{prop:mp_no_pure}
The Matching Pennies game (Example~\ref{ex:mp}) has no pure strategy Nash equilibrium.
\end{proposition}

\begin{proof}
We check all four profiles:
\begin{itemize}[nosep]
    \item $(H,H)$: Player~2 gains by deviating to~$T$ (payoff $-1 \to 1$).
    \item $(H,T)$: Player~1 gains by deviating to~$T$ (payoff $-1 \to 1$).
    \item $(T,H)$: Player~1 gains by deviating to~$H$ (payoff $-1 \to 1$).
    \item $(T,T)$: Player~2 gains by deviating to~$H$ (payoff $-1 \to 1$).
\end{itemize}
Every profile admits a profitable deviation, so no pure NE exists.
\end{proof}

\begin{remark}
The Prisoner's Dilemma has a unique NE at $(D,D)$, the Battle of the Sexes has two
pure NE at $(O,O)$ and $(F,F)$, and the Stag Hunt has two at $(S,S)$ and $(H,H)$.
Matching Pennies has none --- motivating the extension to mixed strategies.
\end{remark}

%% ============================================================
%% SECTION 4: MIXED STRATEGIES
%% ============================================================


\end{document}